\documentclass[main.tex]{subfiles}

\begin{document}
\section{Local Grothendieck-Riemann-Roch theorem}\label{s:GRR}

In this final section we formulate and prove a local, universal form of the Grothendieck-Riemann-Roch theorem for
formal families of complex varieties.
It is a natural generalization of the relationship of the Mumford determinant line bundle defined on the
moduli space of Riemann surfaces and the central charge. 

The set up is as follows.
First, note that $\cJ_d$ is a representation for the dg lie algebra $\lie{witt}_d$
\[
 \sfL \colon \lie{witt}_d \xto{\simeq} \op{Der}( \cJ_d) \to \op{End} \cJ_d .
\]
Pulling back along $\sfL$ is a map in Lie algebra hyper Lie algebra cohomology 
\begin{equation}\label{}
  \sfL^* \colon \HH^2_{Lie}(\op{End} \cJ_d) \to \HH^2_{Lie} (\lie{witt}_d) . 
\end{equation}

More generally, let $\cV$ be any tensor $\lie{witt}_d$-module (recall \ref{TensorMod}). 
Then, we obtain a homomorphism (essentially Lie derivative): 
\begin{equation}\label{}
  \sfL_{V} \colon \lie{witt}_d \to \op{End}(\cV)  
\end{equation}
and hence a map in cohomology 
\begin{equation}\label{}
  \sfL_V^* \colon \HH^2_{Lie}(\op{End} \cV) \to \HH^2_{Lie}(\lie{witt}_d) .
\end{equation}

Assume that $V$ is an $r$-dimensional tensor $\lie{gl}_d$-representation. 
Then, $\rho_V \colon \lie{gl}_d \to \op{End}(V)\simeq \lie{gl}_r$ induces a morphism
\[
\rho_V^*: \op{S}^\bu(\lie{gl}_r^*)^{\lie{gl}_r}\simeq \C[\lie{t}_1,\ldots,\lie{t}_r]\longrightarrow \op{S}^\bu(\lie{gl}_d^*)^{\lie{gl}_d}\simeq \C[y_1,\ldots,y_d].
\]
We define $\op{Ch}_i(\cV) \define \rho_V^*(\lie{t}_i)\in \C[y_1,\ldots,y_d]$. With this in place, the polynomial
\[
\op{Td}\,\op{Ch}(\cV)\in \C[y_1,\ldots,y_d]
\]
is now well-defined.

By \cite[theorem 5.4.16]{FHK} the dg vector space $\op{End} \cV$ is a dg Tate vector space, and therefore its second Lie algebra
cohomology is one-dimensional, with generator we denote 
\begin{equation}\label{}
  \mathfrak{t} \in \HH^2_{Lie}(\op{End}\cV) \simeq \C.
\end{equation}
It is dual to the homology class
\begin{align}\label{eq:universalV}
  \left[z^\1 \wedge \cdots \wedge z^{\ddd} \wedge P \right]  &  \in 
  \HH_2^{Lie}( \op{End} \cV) \\ 
  \<\mathfrak{t}, [z^\1 \wedge \cdots \wedge z^{\ddd} \wedge P] \> & = 1 .
\end{align}
Here, we use that $z^\1, \ldots, z^{\ddd}, P \in \cJ_d$ acts purely through the $\cJ_d$-module structure.
With this setup, we will prove the following local, universal version of the Grothendieck-Riemann-Roch theorem which
was stated as a
conjecture in \cite{KapranovNotes}.

\begin{theorem}\label{thm:grr} 
  Let $V$ be a finite-dimensional $\lie{gl}_d$-representation.
  The composition
\begin{equation}\label{}
  \HH_2^{Lie}(\lie{witt}_d) \xto{(\sfL_V)_*} \HH_2^{Lie}( \op{End} \cV) \xto{\tau} \C ,
\end{equation}
is the following universal class:
  \begin{equation}\label{}
    \op{Td} \op{Ch}(\cV)|_{2d+2} \in \C[\op{ch}_1,\ldots,\op{ch}_d]_{2d+2}
  \end{equation}
  In other words, at the level of cohomology, one has 
  \begin{equation}\label{}
    \sfL^*_V [z^\1 \wedge \cdots \wedge z^{\ddd} \wedge P] = \left[\op{Td} \op{Ch}(\cV)|_{2d+2}\right] \in \HH^2_{Lie}(\lie{witt}_d) .
  \end{equation}
\end{theorem}
%\gui{I think we need to clarify the meaning of Ch for tensor module}

\subsection{Differential operator formulation} 

We sketch our method of proof of theorem \ref{thm:grr} before getting into details.
The main idea is to involve the Jouanolou model of the (dg) algebra of differential operators $\cD_d$ on punctured affine space $\AA^d - \{0\}
$, see \cite[section 1.5]{GWvir2d}.
By construction, $\sfL \colon \lie{witt}_d \ltimes \lie{gl}_r(\cJ_d) \to \op{End} \cJ_d^{\oplus r}$ factors through differential operators
\begin{equation}\label{}
  \lie{witt}_d \ltimes \lie{gl}_r(\cJ_d)\xto{i} \lie{gl}_r(\cD_d) \xto{\til \sfL} \op{End} \cJ_d ^{\oplus r}
\end{equation}
where the first map simply views vector fields acting by first-order differential operators and the second map is the natural inclusion.
From here, we will use theorem \ref{1dImageMatrix} together with the existence of an algebraic trace in the spirit of
\cite{FFS}.
Indeed, in the next section we will define the ``residue trace" on our Jouanolou model for differential operators as a degree two class 
\begin{equation}\label{eq:trD}
  \op{ResTr}_{\cD}^{S^1} \in \HH^2_{Lie}(\lie{gl}_r(\cD_{d})) . 
\end{equation}
Consider the commutative diagram 
\begin{equation}\label{eq:diagram}
  \begin{tikzcd}
    \HH_2^{Lie}(\lie{witt}_d\ltimes \lie{gl}_r(\cJ_d)) \ar[rr,"\sfL_*"] \ar[dr,"i_*"] & & \HH_2^{Lie}(\op{End} \cJ_d^{\oplus r}) \ar[r,"\tau","\cong"'] & \C \\
                                                                     & \HH_2^{Lie}(\lie{gl}_r(\cD_d)) \ar[rr,"\op{ResTr}^{S^1}_{\cD}",
    ""'] \ar[ur,"\til \sfL_*"] & &
                                                                 \C 
                                                                 \end{tikzcd}
                                                               \end{equation}
By definition $\tau([z^1 \wedge z^2 \wedge \cdots \wedge z^d \wedge P]) = 1$
where $\tau$ is the universal Tate class of \cite{FHK}.
We will also show that $\op{ResTr}_{\cD}^{S^1}([z^1 \wedge \cdots \wedge z^d \wedge P]) = 1$.

\begin{theorem}\label{thm:grrD}
The composition 
\begin{equation}\label{}
  \HH_2^{Lie}(\lie{witt}_d\ltimes \lie{gl}_r(\cJ_d)) \xto{i_*} \HH_2^{Lie}(\lie{gl}_r(\cD_d)) \xto{\op{ResTr}_{\cD}^{S^1}} \C .
\end{equation}
is $\op{Td}\cdot \op{Ch}|_{2d+2}$.
\end{theorem}

Theorem \ref{thm:grr} can be reduced to the theorem above. 
Indeed, let $\cV= \cJ_d \otimes V$ be a tensor module. 
The $\lie{witt}_d$-action on $\cV$ factors as a composition
\[
\lie{witt}_d \to \lie{witt}_d \ltimes \lie{gl}_r(\cJ_d)\simeq \lie{witt}_d \ltimes (\operatorname{End}(V)\otimes \cJ_d)\to \op{End}(\cV).
\]
Thus, Theorem \ref{thm:grr} can be deduced from the above theorem.

\subsection{Algebraic residue trace}
As we have already mentioned, this class \eqref{eq:trD} is inspired by the formal algebraic trace map of \cite{FFS}.
We introduce the dg Weyl algebra $\cW_d$ associated to the punctured formal disk, see section \cite[section 1.5]{GWvir2d} for definitions and
notations.
There is a dg algebra isomorphism from $\mathcal{D}_{d}$ to $\mathcal{W}_{d}$ given by the symbol map
\[
\sigma\left(f(z^{\1},...,z^{\dd})F(\partial_{z^{\1}},...,\partial_{z^{\dd}})\right) \define e^{-\frac{1}{2}\sum^d\limits_{\s=1}\partial_{p^{\s}}\partial_{q^{\s}}}(f(q^{\1},...,q^{\dd})F(p^{\1},...,p^{\dd})).
\]
In particular
\[
\sigma\left(T\right)=T-\frac{1}{2}\mathrm{div}(T)
\]
where $T\in \lie{witt}_d \subset \cD_d$ and $\op{div}(\sum_i f^i \del_i) = \sum_i \del_i f^i \in \cJ_d$.

The residue trace is most fundamentally defined at the level of cyclic cohomology.
It is a \textit{non-commutative} version of the cyclic class $\rho \in \HH_\lambda^1(\cJ_d)$ of \cite{FHK}. We follow the presentation in \cite[3]{GLX}, which gives a convenient setting to generalize \cite{FFS} to our case.

\begin{definition}
    Define the residue trace cochain
    \[
      \op{ResTr}^{S^1}_\lambda \in  \overline{C}_{\lambda}^{\bu}(\lie{gl}_r(\mathcal{W}_d))^1
  \]
   of total degree $+1$ by the formula
\begin{multline}
  \mathrm{ResTr}_\lambda^{S^1}\left(O_0\otimes \cdots\otimes O_k\right) =
  \mathrm{Res}_{q} \left({\int_{S^1_{\mathrm{cyc}}[k+1]}\operatorname{tr}_{\lie{gl}_r}\circ\mathbf{Mult}\left(e^{\Pi+D}(O_0\cdot d\theta_0\otimes O_1\otimes\cdots\otimes O_k)\right)}|_{p=0}\right),
\end{multline}
%\begin{multline}
     %mathrm{Tr}^{S^1}\left(O_0\otimes \cdots\otimes O_k\right) \\
%=\mathrm{Res}\left({\int_{\triangle_k}e^{\sum^d\limits_{\s=1}\sum\limits_{i<j}\psi(\theta_i-\theta_j)((\partial_{q^{\s}})_i\otimes(\partial_{p^{\s}})_j-(\partial_{p^{\s}})_i\otimes(\partial_{q^{\s}})_j)}(O_0\otimes D(O_1)\otimes\cdots\otimes D(O_k))}|_{p=0}\right),
%\end{multline}
    where:
    \begin{itemize}
    \item[(i)] the operators $\Pi$ and $D$ are as follows:
      \[
\Pi\define \frac{1}{2}\sum_{i<j}\pi^*_{ij}\SP\cdot \Pi_{ij},\quad \Pi_{ij}:=\sum^d_{\s=1}(\partial_{q^{\s}})_i\otimes(\partial_{p^{\s}})_j-(\partial_{p^{\s}})_i\otimes(\partial_{q^{\s}})_j,
\]
\[
D(O_0\cdot d\theta_0\otimes O_1\otimes\cdots\otimes O_k)=\sum^{k}_{i=1}O_0\cdot d\theta_0\otimes \cdots\otimes D_i(O_i)\otimes\cdots\otimes O_k,
\]
    \[
    D_i(O_i)\define \sum^d_{\s=1}{\partial_{q^{\s}}}(O_i)dq^{\s}\cdot d\theta_i+\sum^d_{\s=1}{\partial_{p^{\s}}}(O_i)dp^{\s}\cdot d\theta_i.
  \]
      \item[(ii)] Let $\mathrm{Cyc}_{S^1}[k+1]$ be the configuration space of $k+1$ anti-clockwise cyclic ordered points on the circle $S^1$. We have a natural identification
      \[
      \mathrm{Cyc}_{S^1}[k+1]\simeq S^1\times \mathring{\Delta}_{k}
      \]
      \[
      (\theta_0,\dots,\theta_k)\rightarrow \{p_0\}\times (\theta_{0,1},\theta_{1,2},\dots,\theta_{k,0}).
      \]
      Here we denote $\Delta_k=\{(\theta_{0,1},\theta_{1,2},\dots,\theta_{k-1,k},\theta_{k,0})\in [0,1]^{k+1}| \theta_{0,1}+\theta_{1,2}+\dots+\theta_{k,0}=1\}$ viewed as a manifold with corners and $\mathring{\Delta}_{k}$ its interior. This allows us to compactify $\mathrm{Cyc}_{S^1}[k+1]$ by $S^1\times\Delta_{k}$, which will be denoted by $S^1_{\mathrm{cyc}}[k+1]$.  
      %$\psi:\R\rightarrow[0,1]$ is the function defined by $\psi(u)=2u+1$ if $-1\leq u <0$ and $\psi(u+1)=\psi(u)$. 
      \item[(iii)] For $0\leq i<j\leq k$, we have the fogetful map
      \[
\pi_{ij}:      S^1_{\mathrm{cyc}}[k+1]\rightarrow       S^1_{\mathrm{cyc}}[2]
      \]
      \[
      (\theta_0,\dots,\theta_k)\rightarrow (\theta_i,\theta_j).
      \]
    Recall  that $S^1_{\mathrm{cyc}}[2]\simeq S^1\times \Delta_1=S^1\times [0,1]$. The function $\SP$ on $S^1_{\mathrm{cyc}}[2]$ is define to be $u-\frac{1}{2}$ where $u\in \Delta_1=[0,1]$.
    
      \item[(iv)] The notation $\mathbf{Mult}\left(-\right)|_{p=0}$ means that we first multiply together the tensor components and then set $p^{\s}=dp^{\s}=0$.
      \item[(v)] $\op{Res}_q$ is the same Jouanolou residue introduced in section \cite[section 2.1]{GWvir2d} except we use the
        variable $q$ instead of $z$.
      \end{itemize}
\end{definition}
From the definition, $\mathrm{ResTr}^{S^1}$ is zero except $k=d$.

\subsection{One-loop result}
Recall the logarithmic Todd polynomial is
\[
\log(\mathrm{Todd}(y))=\frac{y}{2}-\sum_{m\geq 1}\frac{B_{2m}}{2m}y^m
\]
where $B_{2m}$ is $2m$th Bernoulli number.
We proceed to the proof of theorem \ref{thm:grrD}.
\begin{theorem*}[\ref{thm:grrD}]
  One has 
  \begin{equation}\label{}
    \op{ResTr}_{\cD}^{S^1}|_{\lie{witt}_d} = \op{Td}\op{Ch}|_{2d+2}
  \end{equation}
\end{theorem*}

\begin{proof}
 The proof is a straightforward refinement of the argument in \cite{GWWchiral}, which employs a setup similar to that of \cite{FFS} and \cite{GLX}. The main distinction is that we introduce a twist by the symbol map (a secondary difference is that we work with the residue rather than formal geometry). The central idea is to rewrite the exponential in a Wick-type form, which in turn produces a Feynman diagram expansion for the residue trace. Furthermore, because we restrict attention to vector fields instead of general higher-order differential operators, the resulting expansion involves only one-loop graphs.

The term involving $\lie{gl}_r(\cJ_d)\subset\lie{witt}_d\ltimes \lie{gl}_r(\cJ_d)$ cannot be connected by any Wick contraction, since only one-loop diagrams contribute, and contracting the $\lie{gl}_r(\cJ_d)$ part would produce a tree component. Consequently, taking the trace $\operatorname{tr}_{\lie{gl}_r}$ yields the factor $\op{Ch}$. This completes the proof.


\end{proof}

\end{document}

